\documentclass[11pt]{article}
\usepackage[margin=1.05in]{geometry}
\usepackage{amsmath,amssymb,amsthm,mathtools,booktabs,array}
\usepackage[T1]{fontenc}
\usepackage{lmodern,microtype}
\usepackage[colorlinks=true,linkcolor=blue!45!black,urlcolor=blue!45!black,citecolor=blue!45!black]{hyperref}
\usepackage{xcolor}
\newtheorem{theorem}{Theorem}[section]
\newtheorem{lemma}[theorem]{Lemma}
\newtheorem{corollary}[theorem]{Corollary}
\newtheorem{proposition}[theorem]{Proposition}
\theoremstyle{remark}\newtheorem{remark}[theorem]{Remark}
\newcommand{\F}{\mathbb F}
\newcommand{\Z}{\mathbb Z}
\newcommand{\C}{\mathbb C}
\newcommand{\E}{\mathbb E}
\newcommand{\supp}{\operatorname{supp}}
\newcommand{\ord}{\operatorname{ord}}
\newcommand{\one}{\mathbf 1}
\title{Sharp Euler-factor agreement for cubic and quintic Dirichlet character packets\\[3pt]\large A projective coding bound for every prime order, with finite detection under Dirichlet GRH}
\author{Research note for the ProveIt project}
\date{27 November 2026}
\begin{document}
\maketitle
\begin{abstract}
For a group of Dirichlet characters of prime exponent $\ell$ and rank $r$, consider distinct multisets of $k$ characters and the density of primes at which their products of Euler factors disagree. Write $k=\ell^a b+u$ with $\ell^a\le k<\ell^{a+1}$, $1\le b\le\ell-1$, $0\le u<\ell^a$. We prove the lower bound $(\ell-1)(\ell+1-b)/\ell^{a+2}$ when $a<r-1$, and $(\ell-1)/\ell^r$ otherwise. The former bound is attained for $b\in\{1,2,3,4,6,\ell-1\}$ (when admissible), and whenever $b$ divides $\ell-1$ or $\ell+1$. The latter bound is always attained. In particular this gives the exact minimum for every rank and packet size at orders three and five, answering a part of the odd-order question in the merged ProveIt GRH article. The density theorem is unconditional; the analytic transfer of the vendored Dirichlet GRH theorem gives finite distinguishing primes and a uniform error-correcting guarantee at scale $O_w(\min(k,\ell^{r-1})^2\log^2(2Q))$. The method couples cyclotomic divisibility, Newton identities, projective Reed--Muller minimum distance and explicit constructions in the integral group ring.
\end{abstract}

\section{Input and scope}
The repository's vendored readable Lean development is pinned, according to its README, to upstream commit \texttt{72fead102f663d5c7a407}\allowbreak\texttt{964a13dc0cf9d977207} of 25 November 2026 \cite{vendor}. Its \texttt{src/Statements.lean} and \texttt{src/Main.lean} export RH for $\zeta$ and two forms of Dirichlet GRH for all moduli and characters, including imprimitive ones \cite{statements,main}. In particular \texttt{LFunction\_ne\_zero\_of\_}\allowbreak\texttt{half\_lt\_re} excludes zeros in $\Re s>1/2$, with the principal-character pole excepted; \texttt{LFunction\_eq\_zero\_re} puts all zeros in $0<\Re s<1$ on the line. The upstream README reports independent checker and Comparator acceptance with the three standard Lean axioms. This note does not rerun the upstream build or independently audit its proof. We use only these statements, not RH for general Artin, Hecke or automorphic $L$-functions.

The merged article \cite{article} (README: 131 reports, 723 pages) explicitly distinguishes unconditional theorems, consequences of Dirichlet GRH, and reductions whose targets remain open. Its section \texttt{rec.tex}, Theorem \texttt{rec:thm:qmatch}, gives the exact quadratic minimum using a Walsh divisibility argument. Its open questions, item 6 of \texttt{rec:sec:open}, ask for an analogous explicit range for products of characters of odd prime order. We address the minimum (not the full range). We use its prime Gram estimate, \texttt{rec:lem:gram}, for the effective arithmetic part. The new mechanism is to pass from a cyclotomic valuation to a homogeneous polynomial on projective space, then to exploit minimum distance rather than an affine binary code. No claim of priority over all external literature is made.

\section{The sharp density theorem}
Let $Q\ge1$, let $X$ be a subgroup of the group of Dirichlet characters modulo $Q$ isomorphic to $\F_\ell^r$, where $\ell$ is prime and $r\ge1$, and let $V=\widehat X\cong\F_\ell^r$. Fix $\zeta=e^{2\pi i/\ell}$ and choose dual coordinates so $\chi_u(v)=\zeta^{u\cdot v}$. The evaluation map $(\Z/Q\Z)^\times\to V$ is surjective by finite abelian duality. For a prime $p\nmid Q$, denote its image by $[p]$.

A packet $A$ is a multiset of $k\ge1$ members of $X$. Its local polynomial at $v$ is
\[
P_A(v,T)=\prod_{u\in A}(1-\zeta^{u\cdot v}T).
\]
Distinct packets $A,C$ of the same size have a signed difference $c_u=\operatorname{mult}_A(u)-\operatorname{mult}_C(u)$, with $\sum c_u=0$. Write $m=\frac12\sum|c_u|$, the common remaining mass after cancellation, so $1\le m\le k$. Put
\[
H(v)=\sum_u c_u\zeta^{u\cdot v},\qquad h_j(v)=\sum_{u\cdot v=j}c_u\quad (j\in\F_\ell).
\]
The disagreement set $S=\{v:P_A(v,T)\ne P_C(v,T)\}$ equals $\supp H$: indeed, equality of polynomials is equality of the multiplicities of the $\ell$ possible roots, i.e. $h_j=0$ for every $j$. If $H(v)=0$, the minimal polynomial $1+T+\cdots+T^{\ell-1}$ implies that all $h_j$ are equal, and their zero sum makes them zero. Conversely the implication is immediate. Thus $S$ is invariant under nonzero scalar multiplication and excludes $0$. By the prime number theorem in arithmetic progressions, the natural relative density among primes is
\begin{equation}\label{eq:density}
\delta(A,C)=|S|/\ell^r.
\end{equation}
Fourier inversion ensures $H\not\equiv0$. All the following density claims are unconditional.

\begin{theorem}[Projective lower bound and sharp phases]\label{thm:main}
Define $a=\lfloor\log_\ell k\rfloor$ and $b=\lfloor k/\ell^a\rfloor$. For any distinct packets of size $k$,
\begin{equation}\label{eq:main}
\delta(A,C)\ge D_\ell(r,k):=
\begin{cases}
\dfrac{(\ell-1)(\ell+1-b)}{\ell^{a+2}},&a\le r-2,\\[6pt]
\dfrac{\ell-1}{\ell^r},&a\ge r-1.
\end{cases}
\end{equation}
The lower bound is attained in every saturated case $a\ge r-1$. In the nonsaturated case it is attained for all admissible $b$ in $\{1,2,3,4,6,\ell-1\}$ and for $b\mid(\ell-1)$ or $b\mid(\ell+1)$. The extremizers are constructible in every $X$ of the specified rank and may be padded with common factors. The statement also holds with $k$ replaced by $m$ for any given pair.
\end{theorem}

\begin{corollary}[Complete cubic and quintic minima]\label{cor:exact}
For $\ell=3$ and $a<r-1$, the minimum is
\[
\delta_{\min}=
\begin{cases}2/3^{a+1},&3^a\le k<2\cdot3^a,\\
4/3^{a+2},&2\cdot3^a\le k<3^{a+1}.
\end{cases}
\]
For $\ell=5$ and $a<r-1$, it is $4(6-b)/5^{a+2}$ for $b=1,2,3,4$ and $b5^a\le k<(b+1)5^a$. In either case it saturates at $(\ell-1)/\ell^r$ exactly when $k\ge\ell^{r-1}$. The greatest possible agreement density of distinct packets is $1-\delta_{\min}$.
\end{corollary}

\begin{center}
\small
\begin{tabular}{ccll}
\toprule
Order & Packet size & Minimum disagreement & Maximum agreement\\
\midrule
3 & 1 & $2/3$ & $1/3$\\
3 & 2 & $4/9$ & $5/9$\\
3 & 3--5 & $2/9$ & $7/9$\\
3 & 6--8 & $4/27$ & $23/27$\\
3 & 9--17 & $2/27$ & $25/27$\\
5 & 1 & $4/5$ & $1/5$\\
5 & 2 & $16/25$ & $9/25$\\
5 & 3 & $12/25$ & $13/25$\\
5 & 4 & $8/25$ & $17/25$\\
5 & 5--9 & $4/25$ & $21/25$\\
\bottomrule
\end{tabular}
\end{center}
\noindent The table applies for $r\ge3$ (the cubic last row can already saturate at $r=3$). No assumptions on zero simplicity or independence of zero ordinates occur.

\subsection{A cyclotomic mass bound}
Set $R=\Z[\zeta]$, $\pi=1-\zeta$. Recall $R/(\pi)\cong\F_\ell$ and $(\ell)=(\pi)^{\ell-1}$ as ideals; this follows by multiplying $1-\zeta^j$ for $1\le j<\ell$, each associated to $\pi$. Let $v_\pi$ denote the valuation at this prime, with $v_\pi(0)=+\infty$.

\begin{lemma}[Augmentation and valuation]\label{lem:mass}
Let $h=(h_0,\ldots,h_{\ell-1})\in\Z^\ell$ have zero sum and $H_h=\sum h_j\zeta^j\ne0$. If $\sum_{h_j>0}h_j\le k$ and $a,b$ are as in Theorem~\ref{thm:main}, then
\[
1\le v_\pi(H_h)\le a(\ell-1)+b.
\]
\end{lemma}
\begin{proof}
The map from the augmentation lattice $I=\{h:\sum h_j=0\}$ to $\pi R$, $h\mapsto H_h$, is an isomorphism of abelian groups. Its kernel is zero by the minimal polynomial; its image is contained in $\pi R$ and contains every $\pi\zeta^j=\zeta^j-\zeta^{j+1}$ (indices modulo $\ell$), hence all of $\pi R$. In particular
\begin{equation}\label{eq:aug}
H_h\in\ell^d\pi R\quad\Longleftrightarrow\quad h\in\ell^d I.
\end{equation}
If $b=\ell-1$, valuation at least $(a+1)(\ell-1)+1$ would imply by \eqref{eq:aug} that every $h_j$ is divisible by $\ell^{a+1}$, impossible for a nonzero histogram of mass less than $\ell^{a+1}$.

If $b\le\ell-2$, suppose the valuation is at least $a(\ell-1)+b+1$. Then $h=\ell^a h'$ with $h'\in I$ and mass $K\le b$. Furthermore $H_{h'}\in\pi^{b+1}R$. The quotient $R/(\pi^{b+1})$ is $\F_\ell[Y]/(Y^{b+1})$ with $Y=\pi$: it has characteristic $\ell$, and the obvious surjection is an isomorphism by the filtration with successive quotients $\F_\ell$. Expand $\zeta^j=(1-\pi)^j$ there to obtain
\[
\sum_j h'_j\binom{j}{d}=0\pmod\ell\qquad(0\le d\le b).
\]
The binomial polynomials of degrees $0,\ldots,b$ form a basis over $\F_\ell$ since $b<\ell$. Thus the positive and negative parts of $h'$ are two multisets of $K$ elements of $\F_\ell$ with the same first $K$ power sums. Newton identities (all $1,\ldots,K$ invertible) give identical monic root polynomials, hence identical multisets. This implies $h'=0$, a contradiction.
\end{proof}

\subsection{From valuation to a projective code}
\begin{lemma}[Homogeneous reduction]\label{lem:hom}
Let $t=\min_{v:H(v)\ne0}v_\pi(H(v))$. The function $f(v)=\pi^{-t}H(v)\bmod\pi$ is nonzero, vanishes at $0$, has reduced polynomial degree at most $t$, and is represented on $V$ by a homogeneous polynomial of degree exactly $t$ over $\F_\ell$.
\end{lemma}
\begin{proof}
Nonvanishing follows by the choice of $t$, and $H(0)=0$. Use coordinate representatives $v_i\in\{0,\dots,\ell-1\}$ and the binomial interpolation basis $\prod_i\binom{v_i}{j_i}$, with $0\le j_i<\ell$. The coefficient of such a term in $H$ is
\[
A_j=\sum_u c_u\prod_i(\zeta^{u_i}-1)^{j_i},
\]
where a zeroth power is $1$. It is divisible by $\pi^{\sum j_i}$. It is also divisible by $\pi^t$, since the coefficient is an integer linear combination of values of $H$ (finite differences on the grid starting at $0$). After dividing by $\pi^t$ and reducing modulo $\pi$, only $\sum j_i\le t$ can survive; the binomial denominators $j_i!$ are invertible. This proves the degree claim.

For $\lambda\in\F_\ell^\times$, the automorphism $\sigma_\lambda$ sending $\zeta$ to $\zeta^\lambda$ satisfies $\sigma_\lambda(\pi)/\pi\equiv\lambda\pmod\pi$ and acts trivially on $R/(\pi)$. Hence
\begin{equation}\label{eq:scaling}
f(\lambda v)=\lambda^t f(v).
\end{equation}
Uniqueness of the reduced polynomial (degree at most $\ell-1$ in each variable) implies that every monomial occurring has total degree congruent to $t$ modulo $\ell-1$. The constant term is zero. For each monomial of degree $d\le t$, choose a variable already appearing with positive exponent and increase its exponent by $t-d$, a nonnegative multiple of $\ell-1$. This does not alter its values on $\F_\ell$, including at zero, and gives a homogeneous representative of degree $t$.
\end{proof}

We use the classical projective Reed--Muller distance formula of S\o rensen \cite{sorensen,ghorpade}: if a homogeneous degree-$t$ form is nonzero as a function on $\mathbb P^{r-1}(\F_\ell)$, with $1\le t\le(r-1)(\ell-1)$ and $t=d(\ell-1)+e$, $1\le e\le\ell-1$, then the number of projective points of nonvanishing is at least
\begin{equation}\label{eq:prm}
(\ell+1-e)\ell^{r-d-2}.
\end{equation}
For $t\ge(r-1)(\ell-1)+1$ the minimum is one. Here nonvanishing is independent of choice of homogeneous representative; evaluation may use any fixed representative of each projective point. This standard theorem, not a conjecture, is the only coding-theoretic input.

\begin{proof}[Proof of the lower bound in Theorem~\ref{thm:main}]
At any $v$ with $H(v)\ne0$ the positive mass of $h(v)$ is at most $k$. Lemma~\ref{lem:mass} gives $t\le a(\ell-1)+b$. The support of $f$ is contained in that of $H$, and consists of whole nonzero projective lines by \eqref{eq:scaling}. The projective distance formula is nonincreasing with the degree, including at the transition to distance one. If $a\le r-2$, apply \eqref{eq:prm} at the upper bound for $t$ to conclude $|S|\ge(\ell-1)(\ell+1-b)\ell^{r-a-2}$. If $a\ge r-1$, $S$ has at least $\ell-1$ elements since it is nonempty and scalar-invariant. Divide by $\ell^r$ in either case. All steps apply with $m$ in place of $k$.
\end{proof}

\section{Sharp constructions}
We construct integral signed coefficient arrays in the group ring $\Z[X]$. Their positive and negative parts define the two packets; the zero total coefficient condition gives equal size. Extra common entries can pad them to any desired size. If $U\le X$ has dimension $a$ and $D_U=\sum_{u\in U}[u]$, evaluation of $D_U$ on $V$ equals $\ell^a\one_{U^\perp}$. Choose the following base arrays on a complementary subspace to $U$; multiplication by $D_U$ repeats every coefficient in a disjoint coset, so it multiplies mass by $\ell^a$ and disagreement density by $\ell^{-a}$.

\subsection{One to four directions and the six-direction case}
Use independent elements $e_1,e_2\in X/U$, with group ring variables $x=[e_1]$, $y=[e_2]$. For $b=1$ use $1-x$. Its support in the dual is the complement of a hyperplane. For $b=2,3,4,6$, use respectively
\begin{align*}
F_2&=(1-x)(1-y),\\
F_3&=F_2(1-xy),\\
F_4&=F_3(1-xy^{-1}),\\
F_6&=(1-x)(1-y)(1-xy)(1-x^2y)(1-x^3y)(1-x^3y^2).
\end{align*}
For $F_2,F_3,F_4$ the Laurent expansions have $2b$ terms, each $\pm1$, with $b$ of each sign: for instance
\begin{align*}
F_3&=1-x-y+x^2y+xy^2-x^2y^2,\\
F_4&=1+xy^2+x^2y^{-1}+x^3y-y-x^2y^2-xy^{-1}-x^3.
\end{align*}
For completeness, the expansion for $F_6$ is
\begin{align*}
F_6={}&1-x-y+x^4y+x y^2-x^6y^2-x^4y^4\\
 &+x^9y^4+x^6y^5-x^{10}y^5-x^9y^6+x^{10}y^6.
\end{align*}
This last identity is also checked by the accompanying exact polynomial script. Reduction of the exponents modulo $\ell$ cannot increase the positive or negative mass, so the resulting mass is at most $b$. The factor directions are distinct modulo $\ell$ in all admissible cases: for $F_4$ use $\ell\ge5$ and for $F_6$ use $\ell\ge7$, as $b\le\ell-1$. (For $F_6$ the pairwise determinants have nonzero absolute values at most three.) Hence the evaluation is zero on exactly $b$ of the $\ell+1$ lines of the two-dimensional dual, and at the origin. The other lines have $\ell-1$ nonzero points each. Thus the base density is $(\ell-1)(\ell+1-b)/\ell^2$. For $b=1$ this same expression equals $(\ell-1)/\ell$. The mass after lifting is at most $b\ell^a$; if it were smaller, the already-proved bound with the smaller mass would contradict the constructed density. This also shows no cancellation of mass can occur in the indicated cases.

\subsection{The last phase for every prime}
In two coordinates take the difference between the sums of all elements on two distinct one-dimensional character subspaces, e.g.
\[
G=\sum_{j=0}^{\ell-1}x^j-\sum_{j=0}^{\ell-1}y^j.
\]
The common constant cancels, leaving positive and negative mass $\ell-1$. On the dual its values are $\ell(\one_{v_1=0}-\one_{v_2=0})$, nonzero on precisely two projective lines. This is the claimed density for $b=\ell-1$.

\subsection{Divisor phases via split and nonsplit conics}
Suppose $b\mid\ell-1$ and $b<\ell-1$. Take a multiplicative subgroup $B\le\F_\ell^\times$ of order $b$, and $\lambda\in\F_\ell^\times\setminus B$. In the two-dimensional character space take the two disjoint packets
\[
A_0=\{(z,z^{-1}):z\in B\},\qquad C_0=\{(z,z^{-1}):z\in\lambda B\}.
\]
For every $c\in\lambda B$, the involution $z\mapsto c/z$ bijects the two indexing sets, and preserves the projection $z+c/z$. Thus the packets agree on the $b$ distinct projective dual directions $(1,c)$. They agree at the origin as well. The lower bound at size $b$ forces disagreement on at least $\ell+1-b$ projective directions; there are exactly that many remaining. Hence the packets attain it.

If $b\mid\ell+1$ with $b\le\ell-1$, let $E=\F_{\ell^2}$, viewed as a two-dimensional $\F_\ell$-space. Its norm-one group $B_0=\{z:z^{\ell+1}=1\}$ is cyclic of order $\ell+1$. Take a subgroup $B\le B_0$ of order $b$ and a coset $\lambda B\ne B$ in $B_0$. Regard these two sets as packets. The dual directions are the trace functionals $z\mapsto\operatorname{Tr}_{E/\F_\ell}(\alpha z)$ for $\alpha\in E^\times/\F_\ell^\times$. The map $\alpha\mapsto\alpha^{\ell-1}$ identifies this quotient with $B_0$. For each $c\in\lambda B$ choose $\alpha$ with $\alpha^{\ell-1}=c$. Since $z^\ell=z^{-1}$ and $c^\ell=c^{-1}$ on the norm-one group,
\[
\operatorname{Tr}(\alpha c/z)=\alpha^\ell z^\ell+\alpha z=\operatorname{Tr}(\alpha z).
\]
The involution $z\mapsto c/z$ again provides $b$ different agreeing projective directions. As above, the lower bound makes the density exact. Lifting by $D_U$ proves every asserted nonsaturated sharp phase.

Finally, when $k\ge\ell^{r-1}$ choose $U$ of dimension $r-1$ and $e\notin U$. The array $D_U(1-[e])$ has mass $\ell^{r-1}$, and disagreement set the nonzero points on the single line $U^\perp$. This proves all remaining assertions of Theorem~\ref{thm:main} and Corollary~\ref{cor:exact}.

\section{Effective consequences of the vendored Dirichlet GRH}
Fix a nonnegative $C^\infty$ function $w$ supported in $(1/2,1)$, normalized by $\int w=1$. For $F:V\to\C$ put
\[
M_x(F)=\sum_{p\nmid Q}(\log p)w(p/x)F([p]),\qquad
\|F\|_A=\sum_{u\in X}|\widehat F(u)|,
\]
where Fourier coefficients use normalized counting measure. We assume the exported Dirichlet GRH statement only in this section. The prime Gram lemma of \cite[\texttt{rec:lem:gram}]{article} gives an absolute constant $C_w$ depending only on $w$ such that, for $x\ge2$ (enlarging $C_w$ if necessary),
\begin{equation}\label{eq:gram}
|M_x(F)-x\E F|\le C_w\|F\|_A\sqrt{x}\log(2Q).
\end{equation}
For clarity, for a single character this follows by Mellin inversion of $-L'/L$ for its primitive inducing character. The principal pole supplies the main term; GRH places the nontrivial zeros on the line, and the rapid decay of the Mellin transform together with the unit-interval zero count $O(\log(q(2+|t|)))$ bounds their contribution by $O_w(\sqrt{x}\log(2q))$. Shifting to $\Re s=-1/4$ by the functional equation gives $O_w(1+x^{-1/4}\log(2q))$ from the remaining terms. Prime powers cost $O_w(\sqrt{x})$ by Chebyshev, and removing primes dividing $Q$ costs $O_w(\log Q)$. Finally Fourier expansion yields \eqref{eq:gram}. This is a standard explicit-formula argument, also detailed in the cited repository lemma.

\begin{theorem}[Finite distance and adversarial decoding]\label{thm:finite}
Under Dirichlet GRH, for every pair of size-$k$ packets, with $m$ the mass after cancellation, if
\[
\sqrt{x}\ge 4C_w m\log(2Q),
\]
then
\begin{equation}\label{eq:finite}
M_x(\one_S)\ge\frac{\ell-1}{2\ell m}\,x.
\end{equation}
Independently, if $\sqrt{x}\ge4C_w\ell^{r-1}\log(2Q)$, then $M_x(\one_S)\ge (\ell-1)x/(2\ell^r)$ for every distinct pair. In particular, with $K=\min(k,\ell^{r-1})$, at every $x\ge\max(2,16C_w^2 K^2\log^2(2Q))$ there is a distinguishing prime in $(x/2,x)$ for every pair, and the weighted distance is at least $(\ell-1)x/(2\ell K)$. If reported local polynomials of an unknown size-$k$ packet are corrupted at primes of total weight strictly below $(\ell-1)x/(4\ell K)$, it is uniquely recovered by minimum weighted Hamming distance among all such packets.
\end{theorem}
\begin{proof}
Let $J(v)=\sum_j h_j(v)^2$. It is nonnegative and vanishes exactly off $S$, and $J(v)\le2m^2\one_S(v)$ since each of the positive and negative masses of the histogram is at most $m$. Parseval on the cyclic field of exponents, using $\sum h_j=0$, gives
\[
J(v)=\ell^{-1}\sum_{\lambda\in\F_\ell^\times}|H(\lambda v)|^2.
\]
Parseval on $V$ then gives $\E J=(\ell-1)\sum c_u^2/\ell\ge2m(\ell-1)/\ell$ since the coefficients are integers. Also $\|J\|_A\le(\ell-1)(2m)^2/\ell$, by the triangle inequality and submultiplicativity of the Fourier $A$-norm. Equation \eqref{eq:gram} yields $M_x(J)\ge x m(\ell-1)/\ell$ at the stated threshold. Divide by $2m^2$ to get \eqref{eq:finite}.

For the other bound, take any projective line $L$ contained in $S\cup\{0\}$ with $L\setminus\{0\}\subset S$. Its indicator away from zero is $B=\one_L-\one_{\{0\}}$. Subspace indicators have Fourier norm one, and more precisely $\|B\|_A=2(\ell-1)/\ell$ by direct Fourier transformation (characters trivial on $L$ contribute $(\ell-1)/\ell^r$, all others $-1/\ell^r$). Its mean is $(\ell-1)/\ell^r$. Apply \eqref{eq:gram} to $B$ and use $w\ge0$. Choosing the applicable bound according to $K$ gives the uniform distance. The weighted Hamming triangle inequality proves the last assertion, since two candidate packets each within the corruption radius of the same received word would be closer than the minimum distance. This is an information-theoretic decoder, not a polynomial-time algorithm.
\end{proof}

\begin{theorem}[Sharp constructions have uniform relative prime counts]\label{thm:counts}
For the sharp nonsaturated examples with $b\ge1$ and $k_0=b\ell^a$, one has
\begin{equation}\label{eq:relative}
|M_x(\one_S)-xD_\ell(r,k)|\le 2C_w k_0D_\ell(r,k)\sqrt{x}\log(2Q).
\end{equation}
For saturated examples, the same holds with $k_0=\ell^{r-1}$. Thus for every $0<\epsilon<1$ a relative accuracy $\epsilon$ is guaranteed for these extremal pairs throughout $x\ge\max(2,4C_w^2 k_0^2\epsilon^{-2}\log^2(2Q))$.
\end{theorem}
\begin{proof}
In a two-dimensional base the support is precisely the union of $c=\ell+1-b$ projective lines with their origin removed. Write its indicator as $\sum_{i=1}^c\one_{L_i}-c\one_{\{0\}}$. Its Fourier coefficients are $c(\ell-1)/\ell^2$ at zero, $(\ell-c)/\ell^2$ at each nonzero character annihilating one of the $c$ lines, and $-c/\ell^2$ elsewhere. There are $c(\ell-1)$ of the middle type and $(\ell+1-c)(\ell-1)$ of the last type, whence
\[
\|\one_S\|_A=\frac{2bc(\ell-1)}{\ell^2}
\]
for the base. For the lift, the indicator is a product with a subspace indicator in independent coordinates, of Fourier norm one, and the remaining coordinates contribute a constant function of norm one. Thus the norm is unchanged and the mean becomes $\ell^{-a}c(\ell-1)/\ell^2$. Their ratio is $2b\ell^a$. For the one-coordinate case $b=1$ the base indicator is nonzero on $\F_\ell\setminus\{0\}$, of mean $(\ell-1)/\ell$ and Fourier norm $2(\ell-1)/\ell$, so the same ratio follows. The saturated case was computed in Theorem~\ref{thm:finite}. Insert these norms in \eqref{eq:gram}.
\end{proof}

\section{Interpretation and further questions}
The exact minimum is a staircase rather than a smooth $1/k$ bound. In the cubic case it drops alternately by factors $2/3$ and $1/2$ as $k$ crosses $2\cdot3^a$ and $3^{a+1}$. The quintic case has four phases per power of five. The saturation threshold has a direct information-theoretic interpretation: at least $\ell^{r-1}$ factors are needed for two distinct packets to differ on only one projective direction. Such groups occur arithmetically in every rank, for example from $r$ distinct primes congruent to $1\pmod\ell$ and the Chinese remainder theorem.

The cyclotomic valuation argument improves on a simple Fourier uncertainty bound because it uses integrality and the zero total mass, and it identifies precisely where small packets cannot sustain a high order of vanishing. Projective coding then captures the symmetry under changing the primitive $\ell$-th root. The cancellation constructions for $b=3,4,6$ are the rank-two Weyl-denominator patterns (types $A_2,B_2,G_2$); their role is economical mass, not any additional analytic hypothesis. The conic construction extends the sharp cases beyond those patterns. The finite GRH part shows that a sharp density obstruction also controls the weighted prime measurements available in each dyadic interval. It is not a new proof of GRH and does not settle a named classical conjecture.

Several questions remain:
\begin{enumerate}
\item Is \eqref{eq:main} sharp for every prime $\ell$ and every phase $b$? The first phases not covered by the present constructions include $(\ell,b)=(7,5)$. A potential approach is to classify size-$b$ signed projection trades in the affine plane with exactly $b$ vanishing projection directions.
\item Determine the complete high-agreement spectrum, not just its minimum, for cubic and quintic packets. Projective second and higher weights and the integrality of Fourier inversion may constrain which codewords lift to integral histograms of small mass.
\item Classify all extremal packets up to a common character translation, linear coordinate changes, cancellation and padding. The classification of minimum projective Reed--Muller codewords in \cite{ghorpade} is a useful starting point, but support classification alone need not determine the integral coefficients.
\item Give explicit numerical constants for a specified smooth $w$, or a deterministic efficient decoder for the finite prime measurements. The present $C_w$ is effective in principle via the explicit formula, but is not evaluated.
\item Extend to characters of composite exponent. For such orders a single trace no longer determines the local polynomial, and a simultaneous treatment of cyclotomic components is needed.
\end{enumerate}

\paragraph{Reproducibility.} The companion Python program checks the Laurent identities, enumerates all packet pairs for selected small $(\ell,r,k)$, tests the sharp constructions using exact projection histograms (including conic examples), and checks the formula for the Fourier norm of line-union supports using integer arithmetic. These computations are sanity checks, not substitutes for the proofs. The result has not been formalized in Lean. See \texttt{README.md} and \texttt{verification.txt} in the archive.

\begin{thebibliography}{9}
\bibitem{vendor} ProveIt contributors. Vendored grh-lean, README, readable edition at upstream commit \texttt{72fead102f663d5c7a407964a13dc0cf9d977207}. \url{https://github.com/VladimirReshetnikov/ProveIt/tree/main/Analysis/GeneralizedRiemannHypothesis/vendor/grh-lean}.
\bibitem{statements} L. Alp\"oge and contributors. Vendored \texttt{grh/src/Statements.lean}. \url{https://github.com/lsernau/grh-lean/blob/72fead102f663d5c7a407964a13dc0cf9d977207/grh/src/Statements.lean}.
\bibitem{main} L. Alp\"oge and contributors. Vendored \texttt{grh/src/Main.lean}. \url{https://github.com/lsernau/grh-lean/blob/72fead102f663d5c7a407964a13dc0cf9d977207/grh/src/Main.lean}.
\bibitem{article} ProveIt contributors. \emph{Beyond the Riemann hypothesis for Dirichlet $L$-functions: consequences, transfers and limits}, 26 November 2026. Section \texttt{rec.tex}, labels \texttt{rec:thm:qmatch}, \texttt{rec:thm:qmatchfinite}, \texttt{rec:lem:gram}, \texttt{rec:sec:open}. \url{https://github.com/VladimirReshetnikov/ProveIt/tree/main/Analysis/GeneralizedRiemannHypothesis/Article}.
\bibitem{sorensen} A. B. S\o rensen. Projective Reed--Muller codes. \emph{IEEE Transactions on Information Theory} \textbf{37} (1991), 1567--1576.
\bibitem{ghorpade} S. R. Ghorpade and R. Ludhani. On the minimum distance, minimum weight codewords, and the dimension of projective Reed--Muller codes. \emph{Advances in Mathematics of Communications} \textbf{18} (2024), 360--382. \url{https://doi.org/10.3934/amc.2023035}.
\end{thebibliography}
\end{document}
